\documentclass[10pt,a4paper,american]{exam}
\usepackage{../misc/macros/classhandout}

\title{Applied Cryptography - Problem Set 5: End-to-End Encrypted Cloud Storage}
\author{Nadim Kobeissi}
\subject{Problem set covering end-to-end encrypted cloud storage systems, cryptographic attacks on real-world implementations, and authentication protocols.}
\keywords{end-to-end encryption, cloud storage, MEGA, Nextcloud, WhatsApp, OPAQUE, SRP, password authentication, key management, cryptographic attacks}

\begin{document}
\classhandoutheader
\section*{Problem Set 5: End-to-End Encrypted Cloud Storage}

\begin{tcolorbox}[colframe=OliveGreen!30!white,colback=OliveGreen!5!white]
	\textbf{Instructions:} This problem set covers topics in end-to-end encrypted cloud storage from topic 2.4\footnote{\url{https://appliedcryptography.page/slides/\#2-4}} of the course. Submit your solutions as a neatly formatted PDF. You are encouraged to collaborate with classmates in studying the material, but your submitted solutions must be your own work. For proofs, clearly state your assumptions, steps, and conclusions.
\end{tcolorbox}

\section{Understanding E2EE Cloud Storage}
\begin{questions}
	\question[15] \textbf{Threat Models and Key Management}

	\begin{parts}
		\part[7] \textbf{The Mud Puddle Test:}
		You're designing an E2EE cloud storage service for a healthcare provider that must comply with HIPAA regulations.
		\begin{subparts}
			\subpart Explain the ``mud puddle test'' and why it's fundamental to true E2EE systems. What are the implications for password recovery mechanisms?
			\subpart Your CEO insists on having a ``password recovery'' feature. Design two different approaches that balance security with usability, clearly stating the security trade-offs of each.
			\subpart Compare Apple's Advanced Data Protection approach versus a mandatory E2EE approach. Which would you recommend for the healthcare provider and why?
		\end{subparts}

		\part[8] \textbf{Multi-Device Key Management:}
		Design a key management system for an E2EE cloud storage service that supports seamless multi-device access.
		\begin{subparts}
			\subpart Compare password-based key derivation, device pairing, and cloud keychain approaches. What are the security and usability trade-offs of each?
			\subpart How would you handle device revocation when a user loses their phone? Consider both immediate security needs and the user's ability to recover their data.
			\subpart Design a key rotation strategy that maintains backward compatibility while improving forward secrecy. How do you handle files shared with other users during rotation?
		\end{subparts}
	\end{parts}

	\question[10] \textbf{Authentication Protocols}

	\begin{parts}
		\part[10] \textbf{OPAQUE and Password Authentication:}
		You're evaluating authentication mechanisms for a new E2EE cloud storage service.
		\begin{subparts}
			\subpart Compare traditional password hashing (PBKDF2/Argon2) with OPAQUE. For each approach, explain what the server learns about the user's password and what attacks are possible after server compromise.
			\subpart The OPAQUE protocol uses an OPRF. Explain how the 2HashDH OPRF works and why the server cannot perform offline dictionary attacks.
			\subpart Your team suggests using OPAQUE's shared session key for encrypting the backup key. Explain why this would be catastrophic for security and what should be used instead.
		\end{subparts}
	\end{parts}
\end{questions}

\section{Cryptographic Attacks on Real Systems}
\begin{questions}
	\question[18] \textbf{MEGA Attack Analysis}

	\begin{parts}
		\part[10] \textbf{RSA Key Recovery Attack:}
		You're analyzing MEGA's cryptographic design for vulnerabilities.
		\begin{subparts}
			\subpart MEGA uses AES-ECB to encrypt RSA private keys without integrity protection. Explain how this enables the RSA key recovery attack that works in just 2 login attempts.
			\subpart The attack exploits RSA-CRT implementation. Describe how modifying the parameter $u' \neq u = q^{-1} \bmod p$ creates an oracle that leaks information about the prime factor $q$.
			\subpart MEGA attempted to patch these vulnerabilities but introduced new ones. Explain why ad-hoc fixes to broken cryptographic designs are dangerous.
		\end{subparts}

		\part[8] \textbf{Plaintext Recovery:}
		Once the RSA key is recovered, an attacker can compromise the entire system.
		\begin{subparts}
			\subpart Explain how the recovered RSA key enables building an AES-ECB decryption oracle for all encrypted keys in MEGA's system.
			\subpart Design a simpler, more secure key hierarchy that would prevent these cascading failures.
		\end{subparts}
	\end{parts}

	\question[12] \textbf{Nextcloud Vulnerabilities}

	\begin{parts}
		\part[6] \textbf{Authentication Failures:}
		Nextcloud's E2EE implementation has critical authentication flaws.
		\begin{subparts}
			\subpart Explain the ``Key Insertion Attack'' where the server can substitute metadata keys. Why doesn't RSA-OAEP provide authentication?
			\subpart Describe the ``Ghost Key Attack'' involving missing bounds checking and propose fixes for both attacks.
		\end{subparts}

		\part[6] \textbf{AES-GCM IV Reuse:}
		Nextcloud reuses IVs when files are modified, breaking AES-GCM security.
		\begin{subparts}
			\subpart Explain why IV reuse in AES-GCM is catastrophic. What can an attacker learn from $C_1 \oplus C_2$ when the same IV is used?
			\subpart Design an attack that recovers plaintext when you know one version of a file and observe another version encrypted with the same key and IV.
			\subpart Compare the security impact of IV reuse in AES-GCM versus AES-CBC. Why is GCM's authenticated encryption particularly vulnerable?
		\end{subparts}
	\end{parts}
\end{questions}

\section{WhatsApp Encrypted Backups}
\begin{questions}
	\question[10] \textbf{HSM-Based Key Recovery}

	\begin{parts}
		\part[10] \textbf{System Design and OPAQUE:}
		WhatsApp uses Hardware Security Modules (HSMs) for their encrypted backup system.
		\begin{subparts}
			\subpart Explain WhatsApp's two-option approach: password-based recovery via HSM versus 64-digit user-managed keys. What are the security and usability trade-offs?
			\subpart WhatsApp's specification is ambiguous about whether they use OPAQUE's export key or shared session key for $\mathtt{OPAQUE}_K$. Explain why using the shared session key would completely break the security model.
			\subpart Compare WhatsApp's centralized HSM approach with Signal's SVR3 system that distributes trust across multiple enclaves. What are the advantages of Signal's approach?
			\subpart The HSM enforces attempt limits to prevent brute-force attacks. Calculate how many attempts an attacker would need on average to crack a 6-digit PIN versus a 15-character password.
		\end{subparts}
	\end{parts}
\end{questions}

\section{Designing Secure E2EE Cloud Storage}
\begin{questions}
	\question[15] \textbf{Formal Protocol Design}

	\begin{parts}
		\part[8] \textbf{Building a Secure Protocol:}
		Based on the formal treatment in the literature, sketch out components of a secure E2EE cloud storage protocol.
		\begin{subparts}
			\subpart Define a minimal key hierarchy with: root key $K_r$, file encryption keys $K_f$, and sharing keys $K_s$. How does separating these roles prevent MEGA's cascading failure pattern?
			\subpart Sketch an authentication flow using OPAQUE that derives $K_r$ from the export key. What specific property of OPAQUE prevents offline attacks after server compromise?
			\subpart Identify three key security properties your design achieves that MEGA/Nextcloud lack. Give one specific attack each system suffers that your design prevents.
		\end{subparts}

		\part[7] \textbf{Implementation Considerations:}
		Consider practical aspects of deploying your E2EE cloud storage protocol.
		\begin{subparts}
			\subpart How would you handle password changes without re-encrypting all files? Design a key rotation mechanism that maintains security while being efficient.
			\subpart Your protocol should support file versioning and deduplication. How do you implement these features without leaking information to the server?
			\subpart What formal analysis tools (ProVerif, Tamarin, \fstar) would you use to verify your protocol? What security properties would you prove?
		\end{subparts}
	\end{parts}
\end{questions}

\begin{tcolorbox}[colframe=EarthBrown!30!white,colback=EarthBrown!5!white]
	\section*{Bonus Question}
	\begin{questions}
		\bonusquestion[10] \textbf{The Cungadero Cloud Deluxe Audit}

		\begin{minipage}{0.15\textwidth}
			\centering
			\includegraphics[width=\textwidth]{images/cungadero.png}
		\end{minipage}%
		\hspace{0.05\textwidth}%
		\begin{minipage}{0.73\textwidth}
			Spamton G. Spamton, self-proclaimed \textit{``Number 1 Rated Salesman 1997,''} has just launched his revolutionary E2EE cloud storage service: the Cungadero Cloud Deluxe. His marketing materials promise ``\texttt{[[BIG SHOT]] SECURITY AT [[Hyperlink Blocked]] PRICES!}'' and claim the system uses ``\texttt{MILITARY-GRADE [[Kromer]] ENCRYPTION}'' that's ``\texttt{SO SECURE, EVEN I CAN'T SEE YOUR [[Commemorative Plates]]!}''
		\end{minipage}

		You've been hired to conduct a security audit before the service goes live to his first million customers (he's very optimistic). During your initial meeting, Spamton excitedly shows you his ``innovative'' features:

		\begin{itemize}
			\item Password recovery using ``Hardware-Secured Recovery Module'' (the hardware is a drawer locked with a key).
			\item ``Post-quantum key agreement'' using RSA-4096 with $e = 3$ and PKCS\#1 v1.5 padding.
			\item File sharing implements ``Zero-Knowledge URLs'' by XORing the file key with a timestamp before base64 encoding it in the URL.
			\item Keys protected by ``Distributed Storage Architecture'' where each key is split using Shamir's Secret Sharing with threshold k=1, n=1.
			\item Authentication uses ``OPAQUE-Compatible Protocol'' which is actually just PBKDF2 with 1 iteration where the derived key is sent directly to the server over HTTPS for comparison.
			\item ``Perfect Forward Secrecy'' achieved by generating new AES keys for each file block using \texttt{Math.random()} seeded with the file's creation timestamp.
			\item ``Authenticated Encryption'' using AES-CBC where the MAC is computed as SHA-1(encrypt(plaintext)).
		\end{itemize}

		\textbf{Your Task:} Design a comprehensive security audit plan for the Cungadero Cloud Deluxe. Your response should include:
		\begin{parts}
			\part \textbf{Triage Strategy}: Which components would you prioritize examining first, and why? Consider both the likelihood of vulnerabilities and their potential impact.
			\part \textbf{Audit Methodology}: Describe your approach for each major component (authentication, key management, encryption, file sharing). Include specific tests you would run and tools you would use.
			\part \textbf{Diplomatic Reporting}: How would you explain the security issues to Spamton in a way that a ``[[BIG SHOT]]'' salesman would understand? Provide examples of how to translate technical vulnerabilities into business risks.
			\part \textbf{Remediation Roadmap}: Create a prioritized list of fixes that balances security improvements with Spamton's need to launch quickly. Which vulnerabilities are absolutely critical to fix before launch, and which could be addressed in ``version 2.0''?
		\end{parts}

		Your analysis should reference real vulnerability patterns from MEGA, Nextcloud, and other systems discussed in class, explain how formal verification could have prevented Spamton's mistakes, and include a comparison with properly implemented E2EE systems.
	\end{questions}
\end{tcolorbox}

\end{document}
